%% ================================================================================================
%%  Appendix C. The simulation in detail
%%  Code: simulations/python/cultures.py, experiments.py (--check), voting.py (--check), examples.py
%% ================================================================================================
\section{The simulation in detail}\label{app:simulation}

\paragraph{Elections.} Each election has $n = 101$ voters (odd, to limit pairwise ties) and $m = 8$
options. Seven are open proposals: each is made by a voter drawn at random and placed at that voter's
ideal point plus a Gaussian perturbation of standard deviation 0.25 per coordinate; the eighth, the status
quo, is drawn from $\Normal(0, I_2)$ independently of the voters. Utilities are minus the Euclidean
distance between a voter's ideal point and an option.
\begin{itemize}
  \item \emph{Spatial}: ideal points drawn from $\Normal(0, I_2)$.
  \item \emph{Polarised}: each voter belongs to one of two blocs with probabilities 0.55 and 0.45, centred at
  $(-1.2, 0)$ and $(1.2, 0)$, with ideal points drawn around the bloc centre with standard deviation 0.6
  per coordinate.
  \item \emph{Impartial}: utilities drawn independently and uniformly on $[0,1]$; no spatial structure.
  \item \emph{Spatial (1D)}, for Table~\ref{tab:cw_probability} and the checks: as spatial, on a line.
\end{itemize}
The exact draws are in \texttt{cultures.py}.

\paragraph{Vetoes and grades.} Voter $i$ vetoes the options beyond a personal tolerance: distance above
$\rho_i \sim U(1.5, 2.5)$ in the spatial cultures ($U(0.75, 1.25)$ on the line), utility below
$\tau_i \sim U(0.2, 0.5)$ in the impartial culture. Voters veto 38\% of the options on average in the
spatial culture, 40\% in the polarised and 35\% in the impartial culture. The tolerance does not depend
on the menu, so a voter never vetoes an option ranked above one they accept, and clones inherit the
vetoes of the original. Majority Judgment uses grades 0--5: 0 (reject) for a vetoed option; acceptable
options are graded 1--5 by their utility relative to the voter's tolerance.

\paragraph{Measures.} For a winner $w$ and a score $S$ over options, efficiency is
$(S(w) - \bar S)/(\max_x S(x) - \bar S)$, with $\bar S$ the mean over options \citep{Merrill1984}; $S$
is the sum of utilities (utility efficiency) or the 10th percentile of utilities across voters (worst-off
efficiency). Condorcet efficiency is computed on the elections that have a Condorcet winner. Monte Carlo
standard errors of a share $p$ estimated on $R$ elections are $\sqrt{p(1-p)/R}$: at most 1.1 points for
the 2,000 elections of the performance tables, 1.6 points for 1,000 elections, 2.2 points for 500.

\paragraph{Manipulation.} For a winner $w$ and each other option $y$, the coalition is the set of voters
who strictly prefer $y$ to $w$, or the first 10 of them (10\% of the voters; voters are exchangeable, so
the first in index order are a random subset). The coalition tries three joint ballots: (1) $y$ ranked
first and acceptable, $w$ ranked last and vetoed, other rankings and vetoes sincere; (2) as (1), with
every option but $y$ vetoed; (3) $w$ ranked last and vetoed, the rest sincere. The election is
manipulable if one of them elects $y$. Other strategies may succeed where these fail, so the shares are
lower bounds, as in \citet{GreenArmytageTidemanCosman2016}.

\paragraph{Clones.} Three clones of an option are added: each voter ranks them next to the original, in a
random order among themselves, and grades and vetoes them as the original. A clone of the original winner
counts as the original.

\paragraph{Reproducibility and checks.} Every election is drawn from its own seed, so results do not
depend on the number of processes; \texttt{experiments.py} writes all tables in a few minutes on four cores. Three scripts check the code:
\begin{itemize}
  \item \texttt{voting.py --check}: the max-flow test of the proportional veto core against its
  definition; Condorcet consistency, Smith consistency and Pareto optimality of Schulze, Split Cycle and
  Ranked Pairs; the Schulze path strengths against an enumeration of paths; textbook examples; and
  Propositions~\ref{prop:decisive}--\ref{prop:qvc-clones} on random profiles;
  \item \texttt{experiments.py --check}: the probability of a Condorcet winner under the impartial
  culture with 1,001 voters against the limits for three to five options; the median-voter result in one
  dimension \citep{Black1948}; Proposition~\ref{prop:gridlock}(i); and the consistency of the measures;
  \item \texttt{examples.py --check}: the winners stated in Section~\ref{sec:examples}.
\end{itemize}

% Generated by simulations/python/experiments.py -- do not edit by hand
\begin{table}[htbp]\centering\footnotesize
\setlength{\tabcolsep}{4pt}
\caption{The rules compared: impartial culture}\label{tab:perf_impartial}
\sbox{\tabbox}{\begin{tabular}{l>{\raggedleft\arraybackslash}p{0.087\linewidth}>{\raggedleft\arraybackslash}p{0.087\linewidth}>{\raggedleft\arraybackslash}p{0.087\linewidth}>{\raggedleft\arraybackslash}p{0.087\linewidth}>{\raggedleft\arraybackslash}p{0.087\linewidth}>{\raggedleft\arraybackslash}p{0.087\linewidth}>{\raggedleft\arraybackslash}p{0.087\linewidth}}
\toprule
Rule & Condorcet eff. (\%) & Utility eff. (\%) & Worst-off eff. (\%) & Accepted (\%) & Contested (\%) & Deadlock (\%) & Fallback (\%) \\
\midrule
Plurality & 48.0 & 58.8 & 17.9 & 67.6 & 42.0 & 0.0 & 0.0 \\
IRV & 85.0 & 75.4 & 32.6 & 69.0 & 29.8 & 0.0 & 0.0 \\
Borda & 84.5 & 92.7 & 49.8 & 70.2 & 18.4 & 0.0 & 0.0 \\
Copeland & 100.0 & 88.1 & 46.6 & 70.2 & 17.4 & 0.0 & 0.0 \\
Minimax & 100.0 & 87.2 & 44.7 & 70.0 & 19.5 & 0.0 & 0.0 \\
Schulze & 100.0 & 87.3 & 44.4 & 70.0 & 19.3 & 0.0 & 0.0 \\
Ranked Pairs & 100.0 & 84.8 & 43.4 & 69.7 & 22.8 & 0.0 & 0.0 \\
Split Cycle & 100.0 & 87.3 & 44.9 & 70.1 & 17.9 & 0.0 & 0.0 \\
\midrule
Approval & 59.7 & 80.3 & 49.1 & 71.6 & 3.6 & 0.0 & 0.0 \\
Majority Judgment & 61.9 & 85.5 & 34.5 & 69.9 & 18.9 & 0.0 & 0.0 \\
\midrule
Fallback & 59.7 & 80.3 & 49.1 & 71.6 & 3.6 & 0.0 & 0.0 \\
Fallback (impasse) & 12.8 & 0.7 & -2.6 & 65.1 & 65.0 & 100.0 & 0.0 \\
Majoritarian compromise & 67.2 & 80.7 & 38.1 & 70.0 & 19.2 & 0.0 & 0.0 \\
Sequential veto & 18.5 & 16.4 & 9.6 & 66.0 & 56.3 & 0.0 & 0.0 \\
\midrule
Unanimity & 12.8 & 0.7 & -2.6 & 65.1 & 65.0 & 100.0 & 0.0 \\
PVC + Schulze & 84.6 & 80.2 & 52.7 & 69.8 & 22.7 & 0.0 & 0.0 \\
QVC, any veto (1/n) & 59.7 & 80.3 & 49.1 & 71.6 & 3.6 & 0.0 & 100.0 \\
QVC, 10\% & 59.7 & 80.3 & 49.1 & 71.6 & 3.6 & 0.0 & 100.0 \\
QVC, 25\% & 60.1 & 80.3 & 48.9 & 71.6 & 3.6 & 0.0 & 85.4 \\
QVC, 33\% & 86.0 & 84.5 & 44.6 & 70.7 & 3.6 & 0.0 & 3.6 \\
QVC, 50\% & 100.0 & 87.3 & 44.4 & 70.0 & 19.3 & 0.0 & 0.0 \\
\bottomrule
\end{tabular}}
\usebox{\tabbox}
\tabnotes{2,000 elections, 101 voters, 8 options of which one is the status quo. Condorcet efficiency: share of the elections with a Condorcet winner (here 58.0\%) in which the rule elects it. Utility efficiency: total utility of the winner, 0 for the average option and 100 for the best \citep{Merrill1984}. Worst-off efficiency: the same for the utility of the voter at the 10th percentile. Accepted: average share of the voters who do not veto the winner. Contested: share of the elections whose winner is vetoed by at least a third of the voters. Deadlock: the status quo stays because no proposal passes. Fallback: QVC excluded every option and kept the least vetoed. QVC, $x$: Condorcet with a blocking quota $\theta = x$ (Section~\ref{sec:procedure}). Monte Carlo standard errors are at most 1.1 points.}
{Monte Carlo simulation, cultures of \texttt{cultures.py}, rules of \texttt{voting.py}; \texttt{experiments.py}.}
\end{table}


% Generated by simulations/python/experiments.py -- do not edit by hand
\begin{table}[htbp]\centering\footnotesize
\setlength{\tabcolsep}{4pt}
\caption{QVC: the blocking quota}\label{tab:theta_sweep}
\sbox{\tabbox}{\begin{tabular}{l>{\raggedleft\arraybackslash}p{0.077\linewidth}>{\raggedleft\arraybackslash}p{0.077\linewidth}>{\raggedleft\arraybackslash}p{0.077\linewidth}>{\raggedleft\arraybackslash}p{0.077\linewidth}>{\raggedleft\arraybackslash}p{0.077\linewidth}>{\raggedleft\arraybackslash}p{0.077\linewidth}>{\raggedleft\arraybackslash}p{0.077\linewidth}>{\raggedleft\arraybackslash}p{0.077\linewidth}}
\toprule
Quota & Condorcet eff. (\%) & Utility eff. (\%) & Worst-off eff. (\%) & Accepted (\%) & Contested (\%) & Fallback (\%) & Manip. all (\%) & Manip. 10\% (\%) \\
\midrule
Spatial, 1/n (any veto) & 77.8 & 97.5 & 94.5 & 82.1 & 0.8 & 100.0 & 99.0 & 87.3 \\
Spatial, 5\% & 77.8 & 97.5 & 94.5 & 82.1 & 0.8 & 100.0 & 99.0 & 87.3 \\
Spatial, 10\% & 77.9 & 97.5 & 94.5 & 82.1 & 0.8 & 95.1 & 99.0 & 87.3 \\
Spatial, 15\% & 80.1 & 97.7 & 94.4 & 82.1 & 0.8 & 66.5 & 99.0 & 87.0 \\
Spatial, 20\% & 88.5 & 98.5 & 94.3 & 81.9 & 0.8 & 27.6 & 99.0 & 79.0 \\
Spatial, 25\% & 96.1 & 99.3 & 94.7 & 81.6 & 0.8 & 7.8 & 99.0 & 48.7 \\
Spatial, 30\% & 99.2 & 99.5 & 94.6 & 81.5 & 0.8 & 2.3 & 99.0 & 20.3 \\
Spatial, 33\% & 99.5 & 99.6 & 94.8 & 81.5 & 0.8 & 0.8 & 99.0 & 12.3 \\
Spatial, 40\% & 100.0 & 99.5 & 94.6 & 81.5 & 1.1 & 0.2 & 96.3 & 6.7 \\
Spatial, 50\% & 100.0 & 99.5 & 94.6 & 81.5 & 1.1 & 0.0 & 71.7 & 5.7 \\
Spatial, none (Schulze) & 100.0 & 99.5 & 94.6 & 81.5 & 1.1 & 0.0 & 61.3 & 5.7 \\
\midrule
Polarised, 1/n (any veto) & 71.5 & 94.8 & 95.2 & 78.5 & 5.9 & 100.0 & 99.7 & 66.7 \\
Polarised, 5\% & 71.5 & 94.8 & 95.2 & 78.5 & 5.9 & 100.0 & 99.7 & 66.7 \\
Polarised, 10\% & 71.5 & 94.8 & 95.2 & 78.5 & 5.9 & 98.5 & 99.7 & 66.7 \\
Polarised, 15\% & 72.2 & 94.9 & 95.2 & 78.5 & 5.9 & 83.0 & 99.7 & 66.7 \\
Polarised, 20\% & 75.7 & 95.4 & 94.5 & 78.4 & 5.9 & 51.6 & 99.7 & 65.0 \\
Polarised, 25\% & 84.1 & 96.5 & 92.1 & 78.0 & 5.9 & 25.3 & 99.7 & 54.7 \\
Polarised, 30\% & 91.6 & 97.0 & 89.6 & 77.5 & 5.9 & 10.8 & 99.7 & 37.0 \\
Polarised, 33\% & 94.9 & 97.4 & 88.8 & 77.3 & 5.9 & 5.9 & 99.7 & 29.0 \\
Polarised, 40\% & 99.0 & 97.3 & 86.8 & 77.0 & 9.0 & 0.7 & 99.0 & 15.7 \\
Polarised, 50\% & 100.0 & 97.2 & 86.4 & 77.0 & 9.0 & 0.0 & 82.0 & 10.7 \\
Polarised, none (Schulze) & 100.0 & 97.2 & 86.4 & 77.0 & 9.0 & 0.0 & 52.0 & 9.3 \\
\bottomrule
\end{tabular}}
\usebox{\tabbox}
\tabnotes{1,000 elections per culture (300 for manipulability), 101 voters, 8 options. Quota: the share $\theta$ of the voters whose vetoes exclude an option; 1/n, a single veto; none, no veto (the Schulze rule). Measures as in Table~\ref{tab:perf_spatial}; Manip.: manipulable by all the voters who prefer another option, or by at most 10\% of the voters, as in Table~\ref{tab:manipulability}.}
{Monte Carlo simulation, cultures of \texttt{cultures.py}, rules of \texttt{voting.py}; \texttt{experiments.py}.}
\end{table}


% Generated by simulations/python/experiments.py -- do not edit by hand
\begin{table}[htbp]\centering\small
\caption{Deadlock under unanimity as the assembly grows}\label{tab:deadlock_n}
\sbox{\tabbox}{\begin{tabular}{l>{\raggedleft\arraybackslash}p{0.140\linewidth}>{\raggedleft\arraybackslash}p{0.140\linewidth}>{\raggedleft\arraybackslash}p{0.131\linewidth}>{\raggedleft\arraybackslash}p{0.131\linewidth}>{\raggedleft\arraybackslash}p{0.131\linewidth}}
\toprule
Voters & Status quo Pareto-optimal (\%) & Unanimity: deadlock (\%) & QVC 1/n: fallback (\%) & QVC 25\%: fallback (\%) & QVC 33\%: fallback (\%) \\
\midrule
3 & 7.4 & 43.3 & 29.0 & 29.0 & 29.0 \\
5 & 26.7 & 67.4 & 51.7 & 14.5 & 14.5 \\
11 & 50.5 & 85.8 & 81.2 & 24.2 & 7.7 \\
21 & 70.1 & 94.6 & 97.4 & 15.5 & 6.4 \\
51 & 85.7 & 98.4 & 99.9 & 12.8 & 2.5 \\
101 & 92.8 & 99.6 & 100.0 & 8.9 & 0.8 \\
201 & 95.3 & 100.0 & 100.0 & 7.1 & 0.4 \\
\bottomrule
\end{tabular}}
\usebox{\tabbox}
\tabnotes{1,000 elections per row, spatial culture (2D), 8 options including the status quo. Status quo Pareto-optimal: the status quo lies in the convex hull of the voters' ideal points, so that no option whatsoever is preferred to it by every voter (Proposition~\ref{prop:gridlock}). Unanimity: no proposal is preferred to the status quo by every voter, which stays. QVC: share of the elections in which every option is vetoed by the quota and the least-vetoed options are kept; QVC still decides. Monte Carlo standard errors at most 1.6 points.}
{Monte Carlo simulation, cultures of \texttt{cultures.py}, rules of \texttt{voting.py}; \texttt{experiments.py}.}
\end{table}

